\documentclass[12pt,a4paper]{article}
\usepackage[utf8]{vntex, inputenc}
\usepackage{amsmath, amssymb, anyfontsize, array, color,  enumerate, enumitem, fancyhdr, float, graphicx, hyperref, lastpage, listings, mathtools, mips, natbib, pifont, tikz}
\usepackage[open,openlevel=0]{bookmark}
\usepackage[framemethod=tikz]{mdframed}

\usetikzlibrary{calc}
\newcommand\HRule{\rule{12cm}{1pt}}

\pagestyle{fancy}
\fancyhead{} % clear all header fields
\fancyhead[L]{
	\color{blue}
	\begin{tabular}{rl}
		\begin{picture}(25,15)(0,0)
		\put(0,-8){\includegraphics[width=8mm, height=8mm]{BK.png}}
		\end{picture}
		\begin{tabular}{l}
			\textbf{\bf \ttfamily Ho Chi Minh University of Technology}\\
			\textbf{\bf \ttfamily Faculty of Computer Science \& Engineering}
		\end{tabular}
	\end{tabular}
}
\fancyhead[R]{
	\begin{tabular}{l}
		\tiny \bf \\
		\tiny \bf
	\end{tabular}
}
\fancyfoot{} % clear all footer fields
\fancyfoot[L]{\scriptsize \ttfamily Computer Architecture Week 6}
\fancyfoot[R]{\scriptsize \ttfamily Page {\thepage}/\pageref{LastPage}}
\renewcommand{\headrulewidth}{0.3pt}
\renewcommand{\footrulewidth}{0.3pt}

\definecolor{dkgreen}{rgb}{0,0.6,0}
\definecolor{gray}{rgb}{0.5,0.5,0.5}
\definecolor{mauve}{rgb}{0.58,0,0.82}

\lstset{frame=tb,
	language=[MIPS]Assembler,
	aboveskip=3mm,
	belowskip=3mm,
	showstringspaces=false,
	columns=flexible,
	basicstyle={\small\ttfamily},
	numbers=none,
	numberstyle=\tiny\color{gray},
	keywordstyle=\color{blue},
	commentstyle=\color{dkgreen},
	stringstyle=\color{mauve},
	breaklines=true,
	breakatwhitespace=true,
	tabsize=3,
	numbers=left,
	stepnumber=1,
	numbersep=1pt,
	firstnumber=1,
	numberfirstline=true
}

\begin{document}
\begin{titlepage}
  \begin{tikzpicture}[remember picture, overlay]
    \draw[line width = 2pt] ($(current page.north west) + (1in,-1in)$) rectangle ($(current page.south east) + (-1in,1in)$);
  \end{tikzpicture}

  \begin{center}

    % Upper part of the page
    \textbf{\fontsize{12pt}{1pt}\selectfont HO CHI MINH CITY UNIVERSITY OF TECHNOLOGY}\\[0.5cm]
    {\fontsize{13pt}{1pt}\selectfont Faculty of Computer Science \& Engineering}\\[0.5cm]
    \begin{figure}[H]
      \centering
      \includegraphics[width=1.7in,height=1.7in]{BK.png}
    \end{figure}

    % Title
    \HRule\\[0.5cm]
    { \textbf{\fontsize{25pt}{1pt}\selectfont Computer Architecture Week 6}}\\[0.4cm]

    \HRule\\[0.8cm]
    \begin{minipage}{0.545\textwidth}
      \begin{flushleft}
        \textbf{Authors:}\\
        \begin{tabular}{l l}
          Lưu Nguyễn Hoàng Minh & 1952845 \\
          Vũ Anh Nhi            & 1952380 \\
          Nguyễn Chính Khôi     & 1952793 \\
          Nguyễn Hoàng          & 1952255 \\
        \end{tabular}
      \end{flushleft}
    \end{minipage}
    \begin{minipage}{0.4\textwidth}
      \begin{flushright}
        \textbf{Class:}\\
        CO2008\_CC02\\
        \textbf{Lecturer:}\\
        Kiều Đỗ Nguyên Bình\\

      \end{flushright}
    \end{minipage}

    \vfill

    % Bottom of the page
    \vspace{2cm}
    {\large} %{\large \today}
  \end{center}
\end{titlepage}


\pdfbookmark[0]{Question 1}{ques1}
\section*{Question 1:}
Given the following program with MIPS assembly.
\begin{mdframed}[hidealllines=true,backgroundcolor=magenta!10]
  \begin{lstlisting}
    .data
      s_a: .float 1.1
      s_b: .float 2.2

    .text
      la      $t0, s_a
      lwc1    $f0, 0($t0)
      la      $t0, s_b
      lwc1    $f1, 0($t0)
      add.s   $f2, $f1, $f0

      li      $v0, 10
      syscall
  \end{lstlisting}
\end{mdframed}
Students run the program and answer the following question
\begin{enumerate}
  \item Explain the function of lwc1, add.s instructions.
  \item Can we use lw to load value for \$f0 register? Explain.
  \item Is the result that is stored in \$f2 3.3? If it is not, explain.
  \item Analyze the values that are stored in \$f0, \$f1, \$f2 based on IEEE 754 standard.
\end{enumerate}

\pdfbookmark[0]{Question 2}{ques2}
\section*{Question 2:}
Given the following program with MIPS assembly.
\begin{mdframed}[hidealllines=true,backgroundcolor=magenta!10]
  \begin{lstlisting}
    .data
      d_a: .double 1.1
      d_b: .double 2.2

    .text
      la      $t0, d_a
      ldc1    $f0, 0($t0)
      la      $t0, d_b
      ldc1    $f1, 0($t0)
      add.d   $f2, $f1, $f0

      li      $v0, 10
      syscall
  \end{lstlisting}
\end{mdframed}
Students run the program and answer the following question
\begin{enumerate}
  \item List the differences between this program and the program in Question 1.
  \item Is the result that is stored in \$f2 3.3? If it is not, explain.
  \item Analyze the values that are stored in \$f0, \$f1, \$f2 based on IEEE 754 standard.
  \item Correct the program to add 2 double number 1.5 and 2.25.
\end{enumerate}

\pdfbookmark[0]{Question 3}{ques3}
\section*{Question 3:}
Implement a procedure that convert input degree in Fahrenheit to Celsius.
If the degree in Celsius higher than 99.5, print ``Warning'' to terminal, otherwise print ``Safe'' to terminal.

\begin{mdframed}[hidealllines=true,backgroundcolor=magenta!10]
  \begin{lstlisting}
    # TODO
  \end{lstlisting}
\end{mdframed}

\end{document}
